\begin{textsample}{POS Dim 3 – ac – Score 38.00 – ac/ac\_inf\_15.txt}  \label{ex:f3_pos_029}
15. \textbf{BEBÊS} ( ZERO A 1 ANO E 6 \textbf{MESES} ) 15.1.
Campo de Experiências “ \textbf{Escuta}, fala, pensamento e imaginação ” Desde o nascimento, as \textbf{crianças} participam de situações comunicativas cotidianas com as pessoas com as quais interagem.
As primeiras formas de interação do \textbf{bebê} são os movimentos do seu corpo, o olhar, a postura corporal, o sorriso, o choro e outros recursos vocais, que ganham sentido com a interpretação do outro. \textbf{Progressivamente}, as \textbf{crianças} vão ampliando e enriquecendo seu vocabulário e demais recursos de expressão e de compreensão, apropriando -se da língua materna – que se torna, pouco a pouco, seu veículo privilegiado de interação.
Na Educação \textbf{Infantil}, é importante promover experiências nas quais as \textbf{crianças} possam falar e ouvir, potencializando sua participação na cultura oral, pois é na \textbf{escuta} de histórias, na participação em conversas, nas descrições, nas narrativas elaboradas individualmente ou em grupo e nas implicações com as múltiplas linguagens que a \textbf{criança} se constitui ativamente como sujeito singular e pertencente a um grupo social.
Desde cedo, a \textbf{criança} manifesta \textbf{curiosidade} com relação à cultura escrita : ao ouvir e acompanhar a leitura de textos, ao observar os muitos textos que circulam no contexto familiar, comunitário e escolar, ela vai construindo sua concepção de língua escrita, reconhecendo diferentes usos sociais da escrita, dos gêneros, suportes e portadores.
Na Educação \textbf{Infantil}, a imersão na cultura escrita deve partir do que as \textbf{crianças} conhecem e das \textbf{curiosidades} que deixam transparecer.
As experiências com a literatura \textbf{infantil}, propostas pelo educador, mediador entre os textos e as \textbf{crianças}, contribuem para o desenvolvimento do gosto pela leitura, do estímulo à imaginação e da ampliação do conhecimento de mundo.
Além disso, o contato com histórias, contos, fábulas, poemas, cordéis etc. propicia a familiaridade com livros, com diferentes gêneros literários, a diferenciação entre ilustrações e escrita, a aprendizagem da direção da escrita e as formas corretas de manipulação de livros.
Nesse convívio com textos escritos, as \textbf{crianças} vão construindo hipóteses sobre a escrita que se revelam, inicialmente, em rabiscos e garatujas e, à medida que vão conhecendo letras, em escritas espontâneas, não convencionais, mas já indicativas da compreensão da escrita como sistema de representação da língua. 15.2.
Direitos de aprendizagem e desenvolvimento - Conviver com \textbf{crianças} e \textbf{adultos} em situações comunicativas do cotidiano, desenvolvendo \textbf{progressivamente} confiança em participar de experiências que envolvam o pensamento, a imaginação, sentimentos, narrativas e diálogos, quer seja em momentos individuais ou coletivos. - \textbf{Brincar} e se deliciar com histórias e poesias, com jogos, parlendas, rimas, \textbf{brinquedos} cantados, textos de imagens e escritos, ampliando seu repertório das manifestações culturais, favorecendo sua progressiva aprendizagem dos vários gêneros e formas de expressão : gestual, oral, escrita, \textbf{plástica}, \textbf{dramática} e musical. - Explorar \textbf{gestos}, expressões, rimas, imagens, canções, enredos de histórias, textos que compõem o patrimônio cultural da humanidade ( histórias, lendas, mitos, fábulas, poemas, parlendas, trava-línguas, piadas, adivinhas, músicas etc ), atuando como autores no processo, redescobrindo e transformando à sua maneira e de acordo com suas possibilidades, considerando o momento de seu desenvolvimento. - Participar de rodas de conversa, relatos de experiência, leitura e contação de histórias, saraus, dramatizações, apresentações, construção de narrativas, representação de papeis no faz de conta, exploração de materiais impressos e de variedade linguística, comunicando -se com os outros, trocando informações e agindo no mundo, criando assim a possibilidade de compreensão, de organização de ideias e consciência de si mesmo. - Expressar sentimentos, ideias, \textbf{desejos}, necessidades, dúvidas, informações e descobertas, possibilitando o compartilhamento de saberes e conhecimentos entre todos os envolvidos, sendo sua autoria garantida pelo \textbf{gesto}, pela fala, pelo desenho, pela escrita ( convencional ou não convencional ) ou por outra forma de \textbf{registro}. - Conhecer -se como sujeitos de sua cultura, reconhecendo e dando sentido às intenções comunicativas e de \textbf{escuta}, aos \textbf{gestos}, as emissões sonoras, as narrativas de diferentes gêneros linguísticos, aos diferentes suportes e gêneros textuais, valorizando suas origens. 15.2.1.
Objetivo de aprendizagem e desenvolvimento : ( EI01EF01 ) - Reconhecer quando é chamado por seu nome e reconhecer os nomes de pessoas com quem convive.
Aprendizagens específicas - Apropriar -se, gradativamente, do seu nome e de pessoas com quem convive. - Expressar -se ao ouvir a pronúncia do seu nome e de pessoas com quem convive. - \textbf{Interessar} -se por aprender a falar o próprio nome e das pessoas com quem convive.
Sugestões de experiências - \textbf{Brincar} usando seu nome, de \textbf{colegas} e \textbf{adultos} com quem convive. - Manifestar -se quando ouvir seu nome, de seus \textbf{colegas} e \textbf{adultos} com quem convive. - \textbf{Brincar} usando a palavra cantada para se comunicar. - Responder quando chamado pelo nome em situação de comunicação. - Comunicar -se usando o nome de \textbf{colegas} e \textbf{adultos} com quem convive. - Identificar seus pertences e de seus \textbf{colegas}. 15.2.2.
Objetivo de aprendizagem e desenvolvimento : ( EI01EF02 ) - \textbf{Demonstrar} interesse ao ouvir a leitura de poemas e a apresentação de músicas.
Aprendizagens específicas - Desenvolver a \textbf{escuta} e apreciação de poemas e músicas. - Desenvolver postura de respeito e \textbf{escuta} do outro. - Apropriar -se, gradativamente, dos diversos usos da linguagem.
Sugestões de experiências - Ouvir a leitura de textos poéticos. - Interagir, expressar e escutar o outro em momentos de leitura de poemas e apresentação musical. - Participar de \textbf{brincadeiras} cantadas. - Cantar e ouvir os \textbf{colegas} cantarem. - Apreciar músicas de diversos estilos por meio da audição de CDs, DVDs, rádio, mp3, computador, ou por meio de intérpretes que possam ir à instituição ( pais, irmãos, pessoas da comunidade ). - \textbf{Brincar}, com os \textbf{colegas} e \textbf{adultos}, respondendo a comandos por meio de \textbf{gestos}, \textbf{balbucios}, movimentos ou vocalizações. - Participar de apresentações de músicas, dança e teatro. - Vivenciar \textbf{brincadeiras} com os \textbf{colegas} e \textbf{adultos}, envolvendo canções associadas a \textbf{gestos} e movimentos. 15.2.3.
Objetivo de aprendizagem e desenvolvimento : ( EI01EF03 ) - \textbf{Demonstrar} interesse ao ouvir histórias lidas ou \textbf{contadas}, observando ilustrações e os movimentos de leitura do adulto-leitor ( modo de segurar o portador e de \textbf{virar} as páginas ).
Aprendizagens específicas - Desenvolver a \textbf{escuta} de histórias lidas ou \textbf{contadas}. - Desenvolver, \textbf{progressivamente}, postura de leitor. - Ampliar, gradativamente, o vocabulário de palavras fazendo uso delas. - Desenvolver, \textbf{progressivamente}, a capacidade de \textbf{contar} e/ou recontar as histórias ouvidas ou vivenciadas. - Estabelecer relação entre palavras conhecidas e ilustrações diversas. - Conhecer um conjunto de histórias, formando um repertório de histórias preferidas.
Sugestões de experiências - \textbf{Imitar} o comportamento leitor do \textbf{adulto} ao explorar livros. - Vivenciar momentos de contação ou leitura de histórias variadas. - Escolher e \textbf{manipular} livros : segurar o livro, \textbf{virar} as páginas, ver imagens, indicar com o olhar ou com o \textbf{dedo} figuras de interesse. - Envolver -se, junto com o \textbf{adulto} e os \textbf{colegas}, da construção de personagens, utilizando recursos simples. - Apreciar histórias lidas, \textbf{contadas} com fantoches, representadas em encenações, escutadas em áudios. - \textbf{Brincar} \textbf{imitando} as falas dos personagens que mais gostam. - Compartilhar com seus \textbf{colegas} e \textbf{adultos}, seus interesses sobre as histórias ouvidas. - Nomear imagens que lhes chamam a atenção, manifestando suas emoções a partir das histórias por meio de \textbf{gestos}, movimentos e \textbf{balbucios}. 15.2.4.
Objetivo de aprendizagem e desenvolvimento : ( EI01EF04 ) - Reconhecer elementos das ilustrações de histórias, apontando -os, a pedido do adulto-leitor.
Aprendizagens específicas - Compreender, gradualmente, as relações entre ilustrações e histórias. - Apropriar -se, \textbf{progressivamente}, de repertório de narrações conhecidas. - Estabelecer comparação entre os elementos das histórias lidas e as ilustrações. - Desenvolver, gradativamente, a comunicação e a expressão.
Sugestões de experiências - Participar de momentos de leitura de histórias e de textos que apresentem imagens significativas. - \textbf{Manusear} livros com imagens, observando e apontando fotos e figuras. - Nomear personagens ou objetos conhecidos, em ilustrações dos livros. - Expressar -se verbalmente e apontar figuras de objetos e de pessoas que visualiza em situações diversas. 15.2.5.
Objetivo de aprendizagem e desenvolvimento : ( EI01EF05 ) - \textbf{Imitar} as variações de entonação e \textbf{gestos} realizados pelos \textbf{adultos}, ao ler histórias e ao cantar.
Aprendizagens específicas - Interpretar e representar histórias lidas observando a expressão dos \textbf{adultos}. - Desenvolver a sensibilidade e a criatividade a partir da \textbf{escuta} de histórias e músicas. - Apropriar -se, gradativamente, de diversas maneiras de se expressar com clareza. - Construir atitude de liberdade de expressão vocal e corporal.
Sugestões de experiências - Vivenciar e \textbf{imitar} ações como leitor. - \textbf{Dramatizar} histórias, \textbf{imitar} personagens, realizar \textbf{gestos} e movimentos a partir da \textbf{escuta} de músicas. - Escutar, observar, falar e gesticular, interagindo com histórias, contos de repetição, poemas e imagens. - Repetir acalantos, cantigas de roda, poesias e parlendas, \textbf{brincando} com o ritmo, a sonoridade e a conotação das palavras. - \textbf{Brincar} utilizando \textbf{gestos} e movimentos nas situações de leitura de história. - Comunicar -se e usar palavras aprendidas nas histórias escutadas. - \textbf{Brincar} com enredos, objetos ou adereços, tendo como referência histórias conhecidas. - Explorar livros buscando \textbf{contar} suas histórias fazendo uso de diferentes entonações, \textbf{gestos}, expressões ou movimentos corporais. 15.2.6.
Objetivo de aprendizagem e desenvolvimento : ( EI01EF06 ) - Comunicar -se com outras pessoas usando movimentos, \textbf{gestos}, \textbf{balbucios}, fala e outras formas de expressão.
Aprendizagens específicas - Ampliar a comunicação por meio de vocalização, \textbf{balbucios}, \textbf{gestos} ou movimentos. - Desenvolver a fala, por meio da audição de \textbf{sons} e pronúncias de palavras \textbf{emitidas} pelos \textbf{adultos}. - Desenvolver, qualitativamente, a audição de \textbf{sons} da língua e sonoridade das palavras. - Expressar -se oralmente com \textbf{autoconfiança}. - Utilizar, gradativamente, a fala de forma competente, em diferentes contextos.
Sugestões de experiências - Responder com um olhar, \textbf{gestos}, sorriso ou \textbf{balbucio}, para expressar necessidades, \textbf{desejos} e sentimentos. - Pedir e atender pedidos, por meio de \textbf{gestos}, fala ou \textbf{balbucio}. - Participar de diálogos com seus pares e \textbf{adultos}. - \textbf{Brincar} de jogo simbólico, utilizando a oralidade e a linguagem corporal. - \textbf{Brincar} de jogos que explorem a sonoridade das palavras. - \textbf{Brincar} de \textbf{esconder} parte do corpo e ter prazer em encontrar. - Expressar -se por meio de vocalização, \textbf{balbucios}, \textbf{gestos}, movimentos e marcas gráficas algo que deseja e desagrados. - \textbf{Brincar} com \textbf{brinquedos} \textbf{imitando} \textbf{sons} : vrummm, chuá chuá, zimmm e outros \textbf{sons}. 15.2.7.
Objetivo de aprendizagem e desenvolvimento : ( EI01EF07 ) Conhecer e \textbf{manipular} materiais impressos e audiovisuais em diferentes portadores ( livro, revista, gibi, jornal, cartaz, CD, tablet etc. ).
Aprendizagens específicas - \textbf{Interessar} -se por materiais impressos e audiovisuais em diferentes portadores. - Atribuir sentido e significados aos materiais impressos e audiovisuais em diferentes portadores. - Apropriar -se, \textbf{progressivamente}, do uso e da função de alguns recursos tecnológicos e midiáticos. - Desenvolver, \textbf{progressivamente}, estratégias de leitura ao \textbf{manipular} materiais impressos e audiovisuais em diferentes portadores. - Constituir -se, gradativamente, como leitor.
Sugestões de experiências - Realizar leitura de gravuras e fotografias por meio de recursos físicos e virtuais. - \textbf{Brincar} de faz de conta, utilizando recursos impressos e audiovisuais ( livros, revistas, tablet, aparelho telefônico, computador, gravador, máquina \textbf{fotográfica}, filmadora, celular e outros recursos ). - \textbf{Imitar} ações e comportamentos típicos de leitor, exemplos \textbf{virar} a página, apontar as imagens, usar palavras, \textbf{gestos} ou vocalizações na intenção de ler em voz alta o que está escrito. - Participar de gravação audiovisual, para posterior apreciação de suas falas e imagens. 15.2.8.
Objetivo de aprendizagem e desenvolvimento : ( EI01EF08 ) - Participar de situações de \textbf{escuta} de textos em diferentes gêneros textuais ( poemas, fábulas, contos, receitas, quadrinhos, anúncios etc. ).
Aprendizagens específicas - Desenvolver o gosto e o prazer pela leitura. - Construir repertório literário. - Construir familiaridade com o mundo da escrita por meio da \textbf{escuta} de textos. - Desenvolver atitudes positivas em relação à importância e ao valor da escrita na vida social e individual. - Apropriar -se, \textbf{progressivamente}, dos usos e funções sociais da linguagem escrita.
Sugestões de experiências - Escutar histórias, contos, lendas, fábulas, poemas, parlendas, \textbf{contadas} ou lidas por \textbf{adultos} da instituição ou da comunidade. - Escutar leitura de diferentes gêneros em espaços aconchegantes, internos ou externos. - Participar de apresentações de teatro, encenação com fantoches e dramatizações. - Escutar áudio de histórias ou de canções, poemas, parlendas, dentre outros. - Participar da realização de uma receita de algo para comer ou de uma \textbf{tinta} para misturar. - Divertir -se com a \textbf{escuta} de diferentes gêneros textuais, como parlenda, poema, canções, história, receitas, dentre outros. - \textbf{Brincar} de registrar suas histórias preferidas por meio de fotografias, áudios, desenhos e modelagens. 15.2.9.
Objetivo de aprendizagem e desenvolvimento : ( EI01EF09 ) - Conhecer e \textbf{manipular} diferentes instrumentos e suportes de escrita.
Aprendizagens específicas - Familiarizar -se com diferentes instrumentos e suportes de escrita. - Desenvolver o interesse por explorar os instrumentos e suportes de escrita. - Desenvolver, \textbf{progressivamente}, prazer em \textbf{manipular} diferentes instrumentos e suportes de escrita. - Construir autonomia ao \textbf{manusear} os instrumentos e suportes de escrita. - Desenvolver, \textbf{progressivamente}, o \textbf{hábito} de \textbf{cuidado} com os diferentes instrumentos e suportes de escrita. - Desenvolver, gradativamente, a capacidade de selecionar, experimentar, \textbf{manipular}, comparar e descartar os instrumentos e suportes de escrita. - Interagir e compartilhar instrumentos e suportes de escrita.
Sugestões de experiências - \textbf{Manusear} livros e outros portadores de texto. - Explorar livros de materiais diversos ( \textbf{plástico}, tecido, borracha, papel ). - \textbf{Manusear} \textbf{embalagens} de produtos de supermercado, livros variados, livro-brinquedo, livro de imagem, livros com textos, CDs e recursos audiovisuais em espaços de \textbf{brincadeiras} de faz de conta. - Presenciar situações significativas de leitura e escrita. - Escolher e “ ler ”, à sua maneira, livros e outros portadores de texto. - Solicitar a leitura de texto, apontando para um cartaz, mural, revistas, livros ou imagem. - Solicitar a \textbf{escuta} de poemas e canções, apontando para o portador. - Explorar diferentes instrumentos e suportes de escrita em diversas situações. - \textbf{Brincar} com seus pares e \textbf{adultos}, utilizando diferentes instrumentos e suportes de escrita. - Participar da organização e \textbf{cuidado} dos diferentes instrumentos e suportes de escrita. 15.3.
Avaliação - \textbf{acompanhamento} do processo de aprendizagem e desenvolvimento dos \textbf{bebês} : Observação, \textbf{registro} e análise das atitudes dos \textbf{bebês} quanto : - À intenção comunicativa : nomear pessoas ou gravuras que estão vendo em um álbum de retratos ou numa revista ou em ilustração de livros ; responder ao ser chamado pelo próprio nome ; dirigir o olhar ao ouvir a pronúncia do nome de pessoas com quem convive e pronunciar seu próprio nome e de pessoas com quem convive. - Às manifestações expressivas em contextos reais de enunciação : nas narrativas, nas tentativas de falar sua opinião e em outras manifestações. - À interpretação das ideias, motivações e \textbf{desejos} nas situações de comunicação. - Ao interesse pelas \textbf{brincadeiras} cantadas que fazem parte da cultura. - Ao interesse pela \textbf{escuta} e apreciação de histórias, poemas e músicas. - À interação com os materiais impressos e audiovisuais em diferentes portadores. - Às estratégias que usam para \textbf{brincar} de ler, \textbf{manusear} os diferentes portadores de textos, instrumentos e suportes de escrita.

% matched lemmas: acompanhamento, adulto, autoconfiança, balbucio, bebê, brincadeira, brincar, brinquedo, colega, contar, criança, cuidado, curiosidade, dedo, demonstrar, desejo, dramatizar, dramático, embalagem, emitir, esconder, escuta, fotográfico, gesto, hábito, imitar, infantil, interessar, manipular, manusear, mês, plástico, progressivamente, registro, som, tinta, virar
\end{textsample}
